\chapter*{Передмова}
%\addcontentsline{toc}{chapter}{Передмова}

Цей посібник створено для широкого кола читачів: від звичайних користувачів до освітян та волонтерів. Видання поєднує теоретичні основи, аналіз законодавства та практичні кейси, що ілюструють типові сценарії кібершахрайства. 

Матеріал розроблено з урахуванням засад \textbf{Національної стратегії кібербезпеки України} та європейської рамки цифрових компетентностей (\textbf{DigComp}). Основна мета — системне формування навичок безпечної поведінки в мережі без зайвої технічної складності.

\vspace{0.5em}

%\subsubsection*{Застереження та авторське право}
\begingroup
\small
\noindent \textit{Усі історії ґрунтуються на фактичних подіях, проте імена та деталі змінено задля конфіденційності. Ілюстративні матеріали (скріншоти, діалоги) створені автором з навчальною метою.}

\vspace{0.5em}

\noindent \textbf{Авторське право.} Видання захищене Законом України «Про авторське право і суміжні права». Будь-яке відтворення, розповсюдження або використання матеріалів без письмової згоди автора заборонено. Цитування допускається виключно з навчальною або науковою метою за умови обов’язкового зазначення авторства та джерела.
\endgroup

\vspace{0.3 em}

Автор відкритий до пропозицій щодо вдосконалення матеріалу. Якщо ви бажаєте долучитися до проєкту або \textbf{поділитися власною історією}, скористайтеся контактами автора.

\vspace{0.3em}

\begin{center}
    \textit{ Швидка та професійна консультація здатна вберегти вас від фатальних помилок, заощадивши не лише гроші, але й кілометри нервів та безліч згаяного часу.
}
\end{center}



\vfill

\noindent \textbf{Бажаю корисного читання! \\ Нехай ваш цифровий простір буде безпечним.}

\clearpage